\documentclass[12pt,a4paper]{article}

%=========================================================
%                   ENCODAGE ET POLICES
%=========================================================
\usepackage[utf8]{inputenc}
\usepackage[T1]{fontenc}
\usepackage[french]{babel}
\usepackage{lmodern}

%=========================================================
%                     MATHÉMATIQUES
%=========================================================
\usepackage{amsmath,amssymb,amsfonts,mathrsfs}
\usepackage{tkz-tab}
\everymath{\displaystyle}

%=========================================================
%                     MISE EN PAGE
%=========================================================
\usepackage[left=1.5cm,right=1.5cm,top=1.5cm,bottom=1.8cm]{geometry}
\usepackage{multicol}
\usepackage{float}
\usepackage{array,multirow}
\usepackage{enumitem}

%=========================================================
%                GRAPHISMES ET COULEURS
%=========================================================
\usepackage{tikz}
\usetikzlibrary{arrows.meta,calc}
\usepackage{pgfplots}
\pgfplotsset{compat=1.15}

\usepackage[most]{tcolorbox}
\tcbuselibrary{skins,hooks}

%=========================================================
%                  LIENS ET ARRIÈRE-PLAN
%=========================================================
\usepackage[colorlinks=true,linkcolor=blue,urlcolor=blue,citecolor=blue]{hyperref}
\usepackage{eso-pic}

%=========================================================
%                   COULEURS ET STYLES
%=========================================================
\definecolor{myred}{RGB}{180, 40, 40}
\definecolor{myblue}{RGB}{20, 50, 120}
\definecolor{mygreen}{RGB}{30, 120, 40}

%=========================================================
%                       FILIGRANE
%=========================================================
\AddToShipoutPicture{
    \AtTextCenter{
        \makebox[0pt]{
            \rotatebox{45}{
                \textcolor[gray]{0.92}{
                    \fontsize{5cm}{5cm}\selectfont PGB
                }
            }
        }
    }
}

\begin{document}

% En-tête
\begin{center}
    \Large\textbf{\underline{\textcolor{myred}{CHAPITRE : PRIMITIVES D'UNE FONCTION}}}\\[0.2cm]
    \normalsize\textbf{Prof : M. BA} \hfill \textbf{Classe : Terminale L}\\[0.1cm]
    \textbf{Durée : 4 heures} \hfill \textbf{Année scolaire : 2026 -- 2027}
\end{center}
\hrule
\vspace{0.3cm}

%=========================================================
%   I. NOTION DE PRIMITIVE
%=========================================================

\section*{\color{myred}I. Notion de primitive d'une fonction}

\subsection*{\color{myred}1. Définition}

\begin{tcolorbox}[colback=blue!5!white,colframe=myblue,title=\textbf{Définition}]
Soit $f$ une fonction définie sur un intervalle $I$.\\
On appelle \textbf{primitive} de $f$ sur $I$, toute fonction $F$ dérivable sur $I$ telle que pour tout $x \in I$ :
\[
F'(x) = f(x)
\]
\end{tcolorbox}

\paragraph{Exemple}
Soit $f(x) = 2x$ définie sur $\mathbb{R}$.\\
Si on pose $F(x) = x^2$, on remarque que $F'(x) = 2x = f(x)$.\\
La fonction $F(x) = x^2$ est donc une primitive de $f(x) = 2x$ sur $\mathbb{R}$.

\subsection*{\color{myred}2. Ensemble des primitives et condition initiale}

\begin{tcolorbox}[colback=red!5!white,colframe=myred,title=\textbf{Propriétés}]
\begin{enumerate}[label=\arabic*.]
    \item Si $F$ est une primitive de $f$ sur $I$, alors toutes les autres primitives de $f$ sont de la forme :
    \[
    x \mapsto F(x) + k \quad \text{où } k \in \mathbb{R} \text{ est une constante réelle.}
    \]
    \item Il existe une \textbf{unique primitive} $F$ de $f$ qui prend une valeur donnée $y_0$ en un point $x_0$ ($F(x_0) = y_0$).
\end{enumerate}
\end{tcolorbox}

\paragraph{Exemple :}
Pour la fonction $f(x) = 2x$, toutes les primitives sont de la forme $F(x) = x^2 + k$.\\
Trouvons la primitive $F$ telle que $F(1) = 5$ :
\[
F(1) = 5 \iff 1^2 + k = 5 \iff 1 + k = 5 \iff k = 4
\]
La primitive unique cherchée est donc $F(x) = x^2 + 4$.

\vspace{0.3cm}
\hrule
\vspace{0.3cm}

%=========================================================
%   II. TABLEAU DES PRIMITIVES USUELLES
%=========================================================

\section*{\color{myred}II. Primitives des fonctions usuelles}

\begin{center}
\renewcommand{\arraystretch}{1.8}
\begin{tabular}{|c|c|c|}
\hline
\textbf{Fonction } $f(x)$ & \textbf{Une primitive } $F(x)$ & \textbf{Intervalle de validité} \\
\hline
$a \quad (a \in \mathbb{R})$ & $ax$ & $\mathbb{R}$ \\
\hline
$x$ & $\dfrac{1}{2}x^2$ & $\mathbb{R}$ \\
\hline
$x^n \quad (n \in \mathbb{N}^*)$ & $\dfrac{1}{n+1}x^{n+1}$ & $\mathbb{R}$ \\
\hline
$\dfrac{1}{x^2}$ & $-\dfrac{1}{x}$ & $]-\infty \,;\, 0[$ ou $]0 \,;\, +\infty[$ \\
\hline
$\dfrac{1}{\sqrt{x}}$ & $2\sqrt{x}$ & $]0 \,;\, +\infty[$ \\
\hline
$\dfrac{1}{x}$ & $\ln|x|$ & $]-\infty \,;\, 0[$ ou $]0 \,;\, +\infty[$ \\
\hline
\end{tabular}
\end{center}

\paragraph{Règles d'opération :}
Si $F$ est une primitive de $f$ et $G$ est une primitive de $g$, alors :
\begin{itemize}
    \item Une primitive de $f + g$ est $F + G$.
    \item Une primitive de $k \times f$ est $k \times F$ (avec $k \in \mathbb{R}$).
\end{itemize}

\vspace{0.3cm}
\hrule
\vspace{0.3cm}


%=========================================================
%   III. APLICATIONS PAR TYPE DE FONCTIONS (SÉRIE L)
%=========================================================

\section*{\color{myred}III. Calcul de primitives par cas simples}

\subsection*{\color{myred}1. Cas des fonctions polynômes}

Pour trouver la primitive d'un polynôme, on intègre terme à terme en utilisant la formule $\dfrac{x^{n+1}}{n+1}$.

\paragraph{Exemple 1 :} Déterminer une primitive $F$ de $f(x) = 3x^2 - 4x + 5$ sur $\mathbb{R}$.
\paragraph{Résolution :}
\begin{itemize}
    \item La primitive de $3x^2$ est $3 \times \dfrac{x^3}{3} = x^3$.
    \item La primitive de $-4x$ est $-4 \times \dfrac{x^2}{2} = -2x^2$.
    \item La primitive de $5$ est $5x$.
\end{itemize}
Ainsi, une primitive de $f$ est : $\mathbf{F(x) = x^3 - 2x^2 + 5x}$.

\paragraph{Exemple 2 :} Déterminer la primitive $G$ de $g(x) = x^3 + 2x - 1$ sur $\mathbb{R}$ telle que $G(0) = 3$.
\paragraph{Résolution :}
Toutes les primitives de $g$ s'écrivent :
\[
G(x) = \frac{x^4}{4} + x^2 - x + k
\]
Utilisons la condition $G(0) = 3$ :
\[
G(0) = \frac{0^4}{4} + 0^2 - 0 + k = 3 \implies k = 3
\]
D'où : $\mathbf{G(x) = \frac{x^4}{4} + x^2 - x + 3}$.

\end{document}